\chapter{Cohomology of the helicoidal carry: evaluation of cocycles and minimum dimension of cumulative memory}
\label{chap:elevacion-cohomologica-dimension-rev2}
\index{dimension!cohomological rank}
\index{carry!discrete integral}
\index{memory!dimensional elevation}

A periodic coordinate records the visible position; the carry records how many times the trajectory has crossed its boundary. The former returns, whereas the latter accumulates. This distinction turns periodicity into a lift: two instants may share a visible phase yet belong to distinct states because they retain different winding numbers.

In the nonadic alphabet, for \(n\in\mathbb Z_{\geq1}\), let us define
\[
 r(n)=1+((n-1)\bmod9),
 \qquad
 q(n)=\left\lfloor\frac{n-1}{9}\right\rfloor,
\]
and the helical lift
\begin{equation}
 H_9(n)=(r(n),q(n)).
 \label{eq:elevacion-helicoidal-nonadica-rev2}
\end{equation}
Then
\[
 H_9(n+9)=H_9(n)+(0,1):
\]
the phase returns and memory changes. If \(c_k\in\mathbb Z\) is the carry of step \(k\to k+1\), its discrete integral
\begin{equation}
 Q_N=Q_0+\sum_{k=0}^{N-1}c_k
 \label{eq:integral-discreta-carry-rev2}
\end{equation}
constitutes the cumulative coordinate of the lift. Its independence is
algebraic; a later metric realization may represent it by means of
a transverse or orthogonal direction.

\section{Independent cocycles and minimum number of memory coordinates}

Let \(G\) be the directed graph of visible states, and let
\(c_1,\ldots,c_d\in Z^1(G;\mathbb Z)\) be integer cocycles. For a route
\(\gamma\), define
\[
 Q_j(\gamma)=\sum_{e\in\gamma}c_j(e),
 \qquad
 \widetilde\gamma=
 \bigl(q_{\rm vis}(\gamma),Q_1(\gamma),\ldots,Q_d(\gamma)\bigr).
\]
On an open route this evaluation retains the representative of the cocycle
and the endpoints. Indeed, with the convention
\((\mathrm d\phi)(e)=\phi(t(e))-\phi(s(e))\),
\begin{equation}
 Q_j^{\,c_j+\mathrm d\phi}(\gamma)
 =Q_j^{\,c_j}(\gamma)+\phi(t\gamma)-\phi(s\gamma).
 \label{eq:cambio-representante-cociclo-ruta-rev2}
\end{equation}
Therefore, for an open path, the representative and its endpoints are fixed,
and the memory preserves those data. If \(\gamma\) is a cycle, the
boundary term vanishes and \(Q_j(\gamma)\) depends solely on the class
\([c_j]\in H^1(G;\mathbb Z)\).

\begin{theorem}[Cohomological criterion for dimensional rank]
\label{thm:rango-dimensional-cohomologico-rev2}
Suppose that the classes of \(c_1,\ldots,c_d\) generate a free submodule of
rank \(d\) in \(H^1(G;\mathbb Z)\). Let
\[
 Q:Z_1(G;\mathbb Z)\longrightarrow\mathbb Z^d
\]
its evaluation on cycles. If a family of \(m\) additive counters is
obtained by means of a homomorphism \(A:\mathbb Z^d\to\mathbb Z^m\), such
that
\[
 \widehat Q:=A\circ Q:Z_1(G;\mathbb Z)\longrightarrow\mathbb Z^m,
\]
and permits reconstruction of the evaluations, that is, there exists
\(B:\mathbb Z^m\to\mathbb Z^d\) with
\(Q=B\circ\widehat Q\), then
\[
 \boxed{m\geq d}.
\]
Therefore, the lift requires at least \(d\) independent memory coordinates
in addition to the visible coordinate.
\end{theorem}

\begin{proof}
The pairing
\(H^1(G;\mathbb Z)\times H_1(G;\mathbb Z)\to\mathbb Z\) and the independence
of \([c_1],\ldots,[c_d]\) imply
\(\operatorname{rk}\operatorname{im}Q=d\). Since
\(Q=B\circ\widehat Q\),
\[
 d=\operatorname{rk}\operatorname{im}Q
 \leq\operatorname{rk}\operatorname{im}\widehat Q\leq m.
\]
No homomorphism into fewer than \(d\) integer counters admits an
additive inverse on the submodule generated by the evaluations \(Q_j\).
The stronger condition \(B\circ A=I_{\mathbb Z^d}\) is sufficient, but the
bound requires only reconstruction of \(Q\) on the evaluations
performed by the cycles.
\end{proof}

The genealogical dimension here quantifies degrees of freedom necessary to
recover the transport. The theorem does not yet identify that rank with a
physical spatial dimension: it provides the minimum domain on which a
metric realization may act without erasing the independent cocycles.

\section{Nonadic tower and projective conservation of unity}

Let
\[
 \mathcal C_9^{(d)}=(\mathbb Z/9\mathbb Z)^d.
\]
For to mark \(\star\in\mathbb Z/9\mathbb Z\), define the inclusion
\[
 h_{d,\star}:\mathcal C_9^{(d)}\longrightarrow
 \mathcal C_9^{(d+1)},
 \qquad h_{d,\star}(x)=(x,\star),
\]
and the projection
\[
 \pi_{d+1,d}(x_1,\ldots,x_{d+1})=(x_1,\ldots,x_d).
\]

\begin{proposition}[Conservation by Dimensional Elevation]
\label{prop:conservacion-elevacion-dimensional-rev2}
We have
\[
 \pi_{d+1,d}\circ h_{d,\star}=\operatorname{id}_{\mathcal C_9^{(d)}}.
\]
Consequently, the cell of dimension \(d\) reappears exactly as a
coordinate section of the cell of dimension \(d+1\): an affine slice, and a
subgroup when \(\star=0\). The term \emph{codimension-one face} is
reserved for an expressly declared additional cubical realization. For the normalized
uniform measures \(\mu_d\),
\[
 (\pi_{d+1,d})_*\mu_{d+1}=\mu_d.
\]
\end{proposition}

\begin{proof}
The first identity follows by eliminating the newly added final
coordinate. For a point \(x\in\mathcal C_9^{(d)}\), the fiber of
\(\pi_{d+1,d}\) contains nine points, each with mass
\(9^{-(d+1)}\); its total mass is \(9^{-d}=\mu_d(\{x\})\).
\end{proof}

The first equality expresses structural conservation; the second,
normalized conservation. Lifting does not replace the preceding layer: it includes it
as a section and distributes the new coordinate over its fibers. In dimension
three, \(|\mathcal C_9^{(3)}|=729\); its decimal completion and the
disjoint decomposition
\[
 1000=729+243+27+1
\]
remain in the residence of cubic completion, where
interior, faces, edges, and added vertex are distinguished.

\section{From the cumulative coordinate to the metric}

Cocycle lifting and the choice of a metric are successive
operations. The former constructs independent memory coordinates; the
latter declares a quadratic form that measures their variations. In APP, the
central block
\[
 B_c=\begin{pmatrix}7&2\\2&7\end{pmatrix}=9P_++5P_-
\]
produce the reference normalization
\[
 M_c:=\frac{B_c}{\sqrt{45}}\in SO^+(1,1),
\]
proved in \cref{prop:app-trit-espectral-markov-lorentz-rev8}.
On the cubic screen, the physical realization selects the
corresponding Lorentzian signature, and the TPK temporal reader orients its
dynamics. The causal chain is
\[
 \begin{aligned}
 \text{return with carry}
 &\longrightarrow\text{cumulative cocycle}
 \longrightarrow\text{memory rank}\\
 &\longrightarrow\text{quadratic form}
 \longrightarrow\text{realized metric}.
 \end{aligned}
\]

\paragraph{Separation between the internal cohomological range and its spatiotemporal realization.}
The nonadic helix and the integral of carry provide the premises for the cohomological rank criterion and its lower bound. The cell tower and conservation of measure follow from the finite projection. Identifying genealogical rank with a spacetime dimension requires the declared metric functor in Dimensional Mechanics.
